    \section{FFT}
    
    \subsection{Simple FFT}
    Multiplicação de dois polinômios em tempo $O((n+m) \cdot \log(n+m))$
    O polinômio $ a + bx + cx^2 + dx^3$ é representado na forma de vetor como $[a, b, c, d]$. Você passa dois $vector$ e a função $conv$ retorna os coeficientes do polinômio resultante.

    FFT é tudo sobre convolução, a multiplicação de dois polinômios é uma convolução. Convolução pode ser feita com funções contínuas ou discretas, a definição para funções contínuas seria:
    
    Dada duas funções contínuas por partes em $[0,\infty]$ a convolução de $f$ e $g$ é definida pela integral
        \[
            (f*g)(t) = \int_0^t f(r)g(t-r)dr
        \]

    Lib description: $fft(a)$ computes $\hat f(k) = \sum_x a[x] \exp(2\pi i \cdot k x / N)$ for all $k$ ($\exp(2\pi i \cdot kx/N) = e^{2\pi i \cdot kx/N} = w^i$, root of unity). $N$ must be a power of 2.
    Useful for convolution:
    \texttt{conv(a, b) = c}, where $c[x] = \sum a[i] \cdot b[x-i]$.
    For convolution of complex numbers or more than two vectors: FFT, multiply pointwise, divide by n, \texttt{reverse(start+1, end)}, FFT back. Rounding is safe if $(\sum a_i^2 + \sum b_i^2)\log_2{N} < 9\cdot10^{14}$
    (in practice $10^{16}$; higher for random inputs).
    Otherwise, use NTT/FFTMod. Time: $\mathcal{O}(N \log N)$ with $N = |A|+|B|$ ($\tilde 1s$ for $N=2^{22}$)
    \inputwithhash{./solutions/Math/fft.cpp}

    \subsection{FFT with random Mod}

     Lib description: Higher precision FFT, can be used for convolutions modulo arbitrary integers as long as $N\log_2(N)\cdot \text{mod} < 8.6 \cdot 10^{14}$ (in practice $10^{16}$ or higher). Inputs must be in $[0, \text{mod})$. Time: $\mathcal{O}(N \log N)$, where $N = |A|+|B|$ (twice as slow as NTT or FFT).

     \inputwithhash{./solutions/Math/fftmod.cpp}
     
    \subsection{FFT to solve String with Wildcards}
    Given two strings $s$ and $t$, with $t$ having wildcards that can match with any character, returns the number of ocurrence of $t$ in $s$.
    Solution $O(N\log_2(N))$.
    
    \bigskip
    
    (Original Problem: Maraton Nacional de Programacion Colombia 2016 - Wildcards).
    \inputwithhash{./solutions/Math/fft_wild_card.cpp}
    
    \subsection{FFT modular (NTT)}

    Lib description: $ntt(a)$ computes $\hat f(k) = \sum_x a[x] \cdot g^{xk}$ for all $k$, where $g=\text{root}^{(mod-1)/N}$. N must be a power of 2. Useful for convolution modulo specific nice primes of the form $2^a b+1$, where the convolution result has size at most $2^a$. For arbitrary modulo, see FFTMod. \texttt{conv(a, b) = c}, where $c[x] = \sum a[i] \cdot b[x-i]$. For manual convolution: NTT the inputs, multiply pointwise, divide by n, \texttt{reverse(start+1, end)}, NTT back. Inputs must be in $[0, \text{mod})$. Time: $\mathcal{O}(N \log N)$.
    
    To use a modular FFT, also known as NTT, you need a $n$-th root of ur modulus and then define some values in general u can say :
    
    $mod = x*2^k$
    
    $G = 3 $ or $5$
    
    $root = G^{mod/(2^k)}$
    
    $root_1 = root^{-1}$
    
    $root_{pw} = 2^k$
    \inputwithhash{./solutions/Math/ntt.cpp}

    \subsection{FWHT (fft using xor/and/or)}
    FFT but the expoent operation is xor/and/or.

    There are some crazy problems where u need to multiply $N$ polynomials of the form $x^0 + x^i$ and the $i$ can be up to $N$. the main approach is to use some D\&C to multiply those and to keep the polynomials divide in two parts. It is also coded down here. Operations can be modular.
    \inputwithhash{./solutions/Math/FWHT.cpp}
